\section{Experiments}
\label{sec:experiments}

We have three sequential experiments.
Section~\ref{sec:exp1} validates the boundary-regret estimator in synthetic settings with closed-form ground truth.
Section~\ref{sec:exp2} tests whether residual predictive expertise and residual decision value diverge on a real human--AI dataset across reward structures.
Section~\ref{sec:exp3} asks whether an ex ante boundary-regret score allocates scarce human review better than standard baselines under asymmetric rewards.
All three experiments run end-to-end on standard CPU workers.

\subsection{Experiment 1: Estimator and Geometry Validation}
\label{sec:exp1}

\paragraph{Question.}
Does the plug-in estimator recover known boundary regret, and remain zero when a safe action dominates the reachable belief region?

\paragraph{Setup.}
We simulate finite-state Bayesian decision problems with known prior, reward matrix, and channels $P(m\mid y), P(h\mid y)$ ($M \perp H \mid Y$).
For each sampled $(y,m,h)$ we compute the exact posteriors $b_m=P(y\mid m)$ and $b_{m,h}=P(y\mid m,h)$, the model-only action $a_m$, and instance-level boundary regret
\[
  \BR(m,h) \;=\; V(b_{m,h}) - r_{a_m}\cdot b_{m,h}.
\]
The estimator does not see the true channels: it estimates $\hat P(m\mid y), \hat P(h\mid y), \hat p(y)$ from finite data via 5-fold cross-fitting with Laplace smoothing, then computes plug-in $\BRhat$ on held-out instances ($N=10{,}000$ per condition).

\paragraph{Conditions.}
The grid varies $G\in\{2,3,5\}$ decision groups and $K\in\{G,2G\}$ latent states, covering the minimal boundary case ($G=2$), the smallest genuinely multi-action case ($G=3$), and a moderate stress test ($G=5$); $K=2G$ checks that estimation does not degrade when multiple latent states share a decision group.
Three roles: C1 is a positive control (actions reward state groups, human signal crosses boundaries), C2 a weak null (prior deep inside one decision region, weak human channel), and C3 a strong null (highly predictive human signal but a safe action dominates every reachable posterior).
For C3 we set $\tau=(1+U)/2$, where $U$ is the maximum reachable posterior mass on any decision group under the exact channels; $\tau>U$ guarantees no reachable posterior can make a non-safe action optimal while the human signal retains its predictive accuracy.

\paragraph{Results.}
Table~\ref{tab:exp1-core} reports metrics averaged across the six configurations; rank correlation is reported only where $\BR$ varies (in both nulls $\BR\equiv 0$, so we report false-positive rate instead).

\begin{table}[htbp]
\centering
\caption{
    Synthetic validation of $\BRhat$ averaged over $G\in\{2,3,5\}$, $K\in\{G,2G\}$.
    C3's $\Delta\mathrm{LL}$ (0.595) exceeds C1's because the human signal is strongly predictive of group membership, yet no posterior crosses the certified threshold.
    $\Delta\mathrm{LL}$: log-loss improvement from $b_m$ to $b_{m,h}$.
}
\label{tab:exp1-core}
\small
\begin{tabular}{@{}llrrrrr@{}}
\toprule
Condition & Role & $\rho(\BRhat,\BR)$ & FPR & $\%\BR{>}0$ & $\%\BRhat{>}0$ & $\Delta$LL \\
\midrule
C1 & Positive control & 0.901        & 0.000 & 37.9 & 37.9 & 0.408 \\
C2 & Weak null        & \text{n/a}   & 0.000 &  0.0 &  0.0 & 0.019 \\
C3 & Strong null      & \text{n/a}   & 0.000 &  0.0 &  0.0 & \textbf{0.595} \\
\bottomrule
\end{tabular}
\end{table}

The estimator separates predictive value from decision value.
In C1, $\BRhat$ ranks true boundary regret well ($\rho=0.901$) across the 37.9\% of instances where crossings occur (per-configuration details in Appendix~\ref{app:exp1-detail}); under the unshifted cross-fitted estimates, FPR is zero in both nulls.
C3 is the critical case: the human signal achieves the highest log-loss improvement of any condition ($\Delta\mathrm{LL}=0.595$), yet $\BRhat$ assigns zero boundary regret to every instance because no posterior crosses the certified threshold.
High predictive accuracy on task-relevant features does not imply decision value when a safe action dominates the reachable belief region.

A paired binary stress test shifts the log-odds of $\hat b_x$ or $\hat b_{x,h}$ by $\pm1$ while holding the simulated cases, fitting protocol, and all other quantities fixed (Appendix~\ref{app:exp1-miscalibration}).
The resulting FPR is component- and geometry-dependent: shifts in $\hat b_x$ yield 49.1--50.9\% FPR in C1, whereas shifts in $\hat b_{x,h}$ yield 23.0--40.4\% FPR in the C3 strong null; the complementary shifts remain zero in these canonical binary cells.

% ==============================================================

\subsection{Experiment 2: Residual Expertise Does Not Equal Decision Value}
\label{sec:exp2}

\paragraph{Question.}
Does the correlation between residual predictive expertise and boundary regret invert across reward regimes on the same model, readers, and data?

\paragraph{Setup.}
We pair a pretrained DenseNet-121 chest X-ray classifier \citep{Huang2017,Cohen2022} with five board-certified radiologists of the CheXpert reader study \citep{Irvin2019} across five official competition conditions (Atelectasis, Cardiomegaly, Consolidation, Edema, Pleural Effusion), yielding $N=25$ reader--task pairs on 500 frontal radiographs.
For each pair, the ground-truth label $y$ is the majority vote of the \emph{remaining four} readers, which eliminates structural $h$--$y$ contamination and ensures $y \perp h$.
Model beliefs $b_x$ are temperature-scaled \citep{GuoEtAl2017} per pair on a 30\% calibration split; evaluation uses the remaining 70\% (${\approx}350$ images per pair).
The human-augmented posterior $\bpost_{x,h} \approx P(Y=1 \mid x, h)$ is estimated via 5-fold cross-fitted logistic regression on features $[\operatorname{logit}(b_x),\, h,\, h \cdot \operatorname{logit}(b_x)]$.
Spearman correlations between residual predictive expertise (log-loss gain over $b_x$) and $\BRhat$ are reported across the 25 pairs under each reward regime; label protocol details in Appendix~\ref{app:exp2-protocol}.

\paragraph{Reward Regimes.}
We evaluate three reward matrices $r_{a,y}$, each placing the decision threshold $r_{0,0}/(r_{0,0}+r_{1,1})$~\citep{PaukerKassirer1980} in a distinct region of the model belief distribution.
R1 is the identity ($r_{0,0}=r_{1,1}=1$, errors $0$), threshold $1/2$, in the bulk.
R2 ($r_{0,0}=1$, $r_{1,1}=5$, errors $0$) penalises false negatives five times as heavily as false positives, threshold $1/6$, below the bulk.
R3 mirrors R2 ($r_{0,0}=5$, $r_{1,1}=1$), threshold $5/6$, in the sparse upper tail.
R3 is the critical negative control: Theorem~\ref{thm:facet-bridge} predicts that when the threshold sits in the tail, facet crossings become rare, boundary-regret mass collapses to near zero, and residual predictive expertise decouples from decision value regardless of how much the reader improves the model's log-loss.

\begin{table}[htbp]
\centering
\caption{Boundary-regret diagnostics on $N=25$ reader--task pairs: actionable BR mass collapses from R1 to R3 and BR concentrates near the active facet in all regimes.
  $n_{\mathrm{v}}$: pairs where $\BRhat\not\equiv 0$; $\bar\rho_{\mathrm{facet}}$: mean Spearman between facet distance and $\BRhat$.
}
\label{tab:exp2}
\small
\begin{tabular}{@{}llrrrrr@{}}
\toprule
Reward & Role & $\bar\rho_{\mathrm{LL}}$ & $\rho{>}0/n_{\mathrm{v}}$ & $\bar\rho_{\mathrm{facet}}$ & Top-decile zero-$\BRhat$ & $n_{\mathrm{v}}$ \\
\midrule
R1 (thr.\ $1/2$) & Primary result    & $\mathbf{+0.477}$ & $25/25$ & $-0.599$ & $11.2\%$          & $25$ \\
R2 (thr.\ $1/6$) & Intermediate      & $-0.051$          & $15/25$ & $-0.866$ & $28.3\%$          & $25$ \\
R3 (thr.\ $5/6$) & Negative control  & $-0.112$          & $0/7$   & $-0.221$ & $\mathbf{97.4\%}$ & $7$  \\
\bottomrule
\end{tabular}
\end{table}

\paragraph{Reward-Conditional Dissociation.}
The correlation between residual predictive expertise (log-loss gain over the model alone) and boundary regret inverts with the reward (Table~\ref{tab:exp2}): $\bar\rho_{\mathrm{LL}}=+0.477$ under R1 (25/25 pairs positive), $-0.051$ under R2 (15/25), and $-0.112$ under R3 (0/7; sign test $p=0.008$, with the remaining 18 of 25 pairs structurally $\BRhat\equiv 0$).
Same model, same readers, same data; only the reward changes.

The mechanism is geometric.
Under R1, threshold $1/2$ sits in a dense region of model beliefs, so a single reader update routinely crosses a facet and predictive expertise aligns with decision value.
Under R3, the false-positive facet sits in the tail and no posterior update reaches it: only 7 of 25 pairs carry any non-zero $\BRhat$, and $97.4\%$ of top-expertise-decile cases (90th percentile of model--reader disagreement $1-b_x[h]$) have $\BRhat=0$.
The R3 correlation $\bar\rho=-0.112$ on these 7 pairs is directionally consistent with the geometric prediction; the collapse of $n_\mathrm{v}$ from 25 to 7 is the direct confirmation that boundary-regret mass has evacuated.
R2 is intermediate ($28.3\%$ top-decile zero-$\BRhat$).
Facet distance correlates negatively with $\BRhat$ under every reward ($-0.599, -0.866, -0.221$ for R1, R2, R3): boundary regret concentrates near the active facet regardless of how predictive the reader is (Theorem~\ref{thm:facet-bridge}).
R2's stronger concentration follows from the same geometry: with the threshold at $1/6$ nearly all decisions default to treat, so the rare facet crossings come precisely from beliefs already near $1/6$. R1's centrally located threshold produces no such concentration.
We call this \emph{reward-conditional dissociation}: a consequence of reward geometry, not a calibration failure.

Appendix~\ref{app:exp2-robustness} shows the inversion and dissociation hold across in- and off-domain pretrained checkpoints (varying the belief geometry) and readers not involved in ground-truth labelling.
As a theory-confirming negative control, Appendix~\ref{app:exp2-cifar10h} evaluates CIFAR-10H: on highly accurate classifiers whose softmax probabilities rarely straddle a reward facet, $\bar\rho$ collapses to near-zero across all rewards ($+0.015$ under R1 vs.\ $+0.477$ on CheXpert R1), confirming that the dissociation signal depends on the model's beliefs straddling the active facet, exactly as Theorem~\ref{thm:facet-bridge} predicts.

\paragraph{Implication for Predictive Audits.}
A residual-expertise audit in the style of \citet{AlurEtAl2023,AlurRaghavanShah2024} would find substantial predictive value here (mean log-loss gain $0.387$ nats across all 25 pairs) and recommend deploying the readers.
Under R3, $97.4\%$ of those same cases carry zero decision value.
The \emph{estimand shift} from predictive to decisional complementarity is not a refinement; it changes the sign of the deployment recommendation.
If this regime-conditional inversion is theorem-predicted, an ex ante $\BRhat$ score should inherit it: dominating reward-aware geometric proximity where boundary-regret mass is reachable, converging with it where mass is absent.
Section~\ref{sec:exp3} tests this directly.

% ==============================================================

\subsection{Experiment 3: Review Allocation under Scarce Human Input}
\label{sec:exp3}

\paragraph{Question.}
Given a fixed human-review budget, can an ex ante boundary-regret score identify the cases where reader input will improve realized decision utility, rather than improve prediction or flag model uncertainty?
The prediction is regime-conditional: it should beat reward-aware but human-blind boundary proximity when reader updates can cross reward-relevant decision boundaries, and converge to it when such crossings are sparse or absent.


\paragraph{Setup.}
We use the same 25 reader--task pairs as in Section~\ref{sec:exp2}.
For each case, the model-only belief $b_x$ induces a model-only action $a_x$.
For each review budget $q\in\{5\%,\,10\%,\,20\%,\,50\%\}$, a policy ranks cases using only pre-review information.
Reviewed cases observe the reader response $h$ and take the human-updated action $\argmax_a r_a\cdot\bpost_{x,h}$; unreviewed cases retain $a_x$.
Utility gain is the held-out realized reward improvement over this no-review baseline.
All learned policy scores are 5-fold cross-fitted on the same fold structure used to estimate $\bpost_{x,h}$, so review priority is always computed out of fold.

\paragraph{Policies.}
The proposed policy, $\BRhat$, is an ex ante boundary-regret score: it predicts possible reader responses from $x$, predicts the corresponding human-updated posterior, and scores a case by the expected decision value of that update under the current reward.
We compare it against six baselines.
\textbf{Margin} is a reward-aware distance to the active decision facet, our primary comparison baseline: it captures reward geometry without estimating the human signal $P(h\mid x)$.
\textbf{Residual} is a residual-expertise score, the deployable analogue of a predictive audit \citep{AlurEtAl2023}.
\textbf{Entropy} is reward-blind model uncertainty.
These policies span two axes: Margin is reward-aware but human-blind, Residual is human-aware but reward-blind, Entropy is blind to both, and $\BRhat$ is both reward-aware and human-aware.
\textbf{L2D} is the learned full-coverage competitor: like $\BRhat$ it is both reward-aware and human-aware, but reaches that coverage end-to-end through a cross-fitted deferral classifier with reward-sensitive labels $\Ind[R[a_{x,h},y]>R[a_x,y]]$, in the spirit of \citep{MozannarSontag2020,MaoMohri2023}, rather than through the reward's geometry and the reader-channel posterior.
\textbf{Random} is the uninformed baseline.
\textbf{Oracle} is the hindsight upper bound using realized labels.

\paragraph{Regime-Conditional Allocation.}
We evaluate all policies on the same 25 reader--task pairs under each reward matrix.
Table~\ref{tab:exp3} tests whether the dissociation in Section~\ref{sec:exp2} becomes actionable under a scarce review budget.
It does, but only where reward geometry leaves actionable mass reachable: $\BRhat$ beats Margin at every R1 budget, leads or ties across most R2 budgets, and has no systematic advantage in the R3 negative control.
The paired wins, losses, and ties isolate when the human-aware term adds value beyond reward geometry.

\begin{table}[htbp]
\centering
\caption{Ex ante review allocation under scarce human input.
  Entries report mean held-out utility gain over no review on $N=25$ reader--task pairs.
  $\BRhat$ beats reward-aware Margin where actionable mass is dense (R1), leads or ties where mass is sparse (R2), and has no systematic edge where mass is absent (R3, negative control).
  W/L/T gives paired wins/losses/ties of $\BRhat$ against Margin; bold marks the best deployable policy in each row (Oracle excluded).
% Under R1, Margin and Entropy induce the same ranking, so their gains coincide.
}
\label{tab:exp3}
\scriptsize
\setlength{\tabcolsep}{3pt}
\begin{tabular}{@{}ll rrrrr@{\hspace{5pt}}|@{\hspace{5pt}}rr@{\hspace{5pt}}|@{\hspace{5pt}}r@{}}
\toprule
Reward & $q$ & $\BRhat$ & Margin & Entropy & Residual & L2D & Random & Oracle & W/L/T \\
\midrule
R1 & $5\%$  & $\mathbf{+0.0459}$ & $+0.0395$ & $+0.0395$ & $+0.0209$ & $+0.0152$ & $+0.0184$ & $+0.0514$ & $20/2/3$ \\
R1 & $10\%$ & $\mathbf{+0.0883}$ & $+0.0743$ & $+0.0743$ & $+0.0455$ & $+0.0289$ & $+0.0392$ & $+0.1000$ & $24/0/1$ \\
R1 & $20\%$ & $\mathbf{+0.1638}$ & $+0.1438$ & $+0.1438$ & $+0.0915$ & $+0.0697$ & $+0.0761$ & $+0.2000$ & $22/1/2$ \\
R1 & $50\%$ & $\mathbf{+0.3325}$ & $+0.2797$ & $+0.2797$ & $+0.2544$ & $+0.2485$ & $+0.2022$ & $+0.4495$ & $25/0/0$ \\
\midrule
R2 & $5\%$  & $\mathbf{+0.0514}$ & $+0.0490$ & $+0.0429$ & $+0.0330$ & $+0.0477$ & $+0.0313$ & $+0.0514$ & $3/0/22$ \\
R2 & $10\%$ & $\mathbf{+0.0999}$ & $+0.0965$ & $+0.0768$ & $+0.0610$ & $+0.0946$ & $+0.0534$ & $+0.1000$ & $3/0/22$ \\
R2 & $20\%$ & $\mathbf{+0.1914}$ & $+0.1889$ & $+0.1486$ & $+0.1286$ & $+0.1879$ & $+0.1089$ & $+0.2000$ & $5/2/18$ \\
R2 & $50\%$ & $+0.4113$ & $\mathbf{+0.4161}$ & $+0.2955$ & $+0.3069$ & $+0.4134$ & $+0.2664$ & $+0.4865$ & $7/9/9$ \\
\midrule
R3 & $5\%$  & $\mathbf{+0.0029}$ & $+0.0027$ & $+0.0000$ & $+0.0000$ & $-0.0008$ & $+0.0006$ & $+0.0072$ & $1/4/20$ \\
R3 & $10\%$ & $+0.0025$ & $\mathbf{+0.0031}$ & $+0.0000$ & $+0.0000$ & $+0.0008$ & $+0.0007$ & $+0.0072$ & $0/4/21$ \\
R3 & $20\%$ & $+0.0025$ & $\mathbf{+0.0029}$ & $+0.0000$ & $+0.0000$ & $+0.0017$ & $+0.0006$ & $+0.0072$ & $1/4/20$ \\
R3 & $50\%$ & $+0.0027$ & $\mathbf{+0.0029}$ & $+0.0000$ & $+0.0000$ & $+0.0023$ & $+0.0021$ & $+0.0072$ & $1/2/22$ \\
\bottomrule
\end{tabular}
\end{table}

\paragraph{Dense actionable mass (R1).}
$\BRhat$ wins decisively at every budget (sign-test $p<0.001$); decisive-win rates over Margin are $91\%$, $100\%$, $96\%$, $100\%$ (tallies $20/2/3$, $24/0/1$, $22/1/2$, $25/0/0$).
With actionable mass dense in the belief distribution, the reader channel adds signal beyond reward-aware geometric proximity: Margin finds cases near the active facet; $\BRhat$ identifies which of those the reader can move profitably.
L2D matches $\BRhat$'s nominal coverage but falls below Random at $q\!\le\!20\%$ ($+0.0152$ vs.\ $+0.0184$ at $q\!=\!5\%$): the structured decomposition is what makes that coverage usable.

\paragraph{Sparse actionable mass (R2).}
The low threshold places most cases on the positive action; sparse boundary-regret mass concentrates near the active facet (exactly where Margin sorts), making Margin a strong baseline. All deployable reward-aware policies obtain much larger absolute gains than under R3.
$\BRhat$ leads or ties Margin at every budget up to $q=20\%$ with no paired losses ($3/0/22$, $3/0/22$, $5/2/18$); the $q=50\%$ row is essentially a tie ($7/9/9$, mean difference $0.005$).
The human-aware term adds value where actionable mass exists and is harmless where mass is sparse.

\paragraph{No actionable mass (R3).}
The active facet sits in the sparse upper tail of model beliefs; eighteen of twenty-five pairs have $\BRhat\equiv 0$ across the evaluation set.
Allocation gains are correspondingly tiny: $\BRhat$ and Margin differ by less than $0.0006$ at every budget, with only three to five decisive pairs out of 25.
This is the no-actionable-mass prediction of Theorem~\ref{thm:facet-bridge}, not an allocation failure.

\paragraph{Human Channel versus Geometric Proximity.}
Reward-aware distance to the active facet is a strong baseline: it already identifies cases near the decision boundary.
The decisive-win rates over Margin under R1 show that $\BRhat$'s estimate of $P(h\mid x)$ and $P(Y\mid x,h)$ add reliable signal on top of that geometric distance where actionable mass is dense.
Where boundary-regret mass is sparse (R2) or absent (R3), $\BRhat$ and Margin converge, both outcomes predicted by Theorem~\ref{thm:facet-bridge}.
Residual and Entropy clarify why both axes matter: human information without reward geometry, or model uncertainty without either reward geometry or the reader channel, allocates review less efficiently.
L2D makes the opposite case: nominal coverage of both axes is not enough on its own; the structured decomposition is what makes that coverage usable at this sample size.
The reward-aware/human-aware policy adds value exactly when both components are needed; it does not manufacture decision value when the reachable belief region carries none.
